\subsection{Complete S-class towers and the first ordinary class groups}

\begin{statement}{E.86}{Proved}{All S-class norms, the ordinary class groups at the second level, and the exact unit-norm defects}
Keep the literal roots of E.76 and E.79:
\[
 K=\mathbb Q(\zeta),\quad \xi^3=\zeta,\quad
 \beta_d^{3^d}=17\zeta,\quad \beta_{d+1}^3=\beta_d,
 \quad L_d=K(\beta_d),\quad B_d=L_d(\xi).
 \tag{E86.1}
\]
Let $S$ mean all primes above $3,17$ in the indicated field.
All class groups below are their three-primary parts. Put
\[
 A_d=\operatorname{Cl}_{S}(B_d)\{3\},\qquad
 C_d^S=\operatorname{Cl}_{S}(L_d)\{3\},\qquad d\ge1.
 \tag{E86.2}
\]
These entire towers, including their actual ideal maps, are
\[
 \begin{gathered}
 A_d=\mathbb Z/3,\qquad C_d^S=(\mathbb Z/3)^2,\\
 N_{B_{d+1}/B_d}=1,\quad j_{B_{d+1}/B_d}=0,
 \qquad N_{L_{d+1}/L_d}=I_2,\quad j_{L_{d+1}/L_d}=0.
 \end{gathered}
 \tag{E86.3}
\]
Here are explicit compatible bases for these formulas. Let
$\mathfrak p_d$ be the degree-one prime of $B_d$ above $73$
with $\xi=55$ and $\beta_d=22$ for odd $d$, $63$ for even
$d$, as in E79.10. Its three-primary S-class is the displayed
generator of $A_d$. In $L_d$, let
\[
 \begin{gathered}
 J_d=(73,\beta_d-r_d),\qquad
 P_d=(7,\beta_d+1),\\
 a_d=[J_d/P_d]_3,\quad b_d=[J_d]_3,\qquad
 r_d=\begin{cases}22&d\text{ odd},\\63&d\text{ even}.
 \end{cases}
 \end{gathered}
 \tag{E86.4}
\]
The ideals mean the kernels of the indicated residue maps on
the maximal orders. Their degree is one and their adjacent
norms are their predecessors. The pair $(a_d,b_d)$ is the
specified S-class basis. The remaining maps at each level are
\[
 \begin{aligned}
 N_{B_d/L_d}:\mathbb Z/3&\longrightarrow(\mathbb Z/3)^2,
                    &z&\longmapsto(0,z),\\
 j_{B_d/L_d}&=0,&
 \operatorname{Art}_{B_d/L_d}(x a_d+y b_d)&=x\pmod3.
 \end{aligned}
 \tag{E86.5}
\]
Arithmetic Artin is used, and $1\in\mathbb Z/3$ sends
$\xi$ to $\zeta\xi$. Thus the image of norm in (E86.5) is
exactly the kernel of this Artin map.

The first ordinary class groups are completely evaluated:
\[
 \operatorname{Cl}(B_2)\{3\}=C_9\oplus C_3^2,
 \qquad \operatorname{Cl}(L_2)\{3\}=C_9\oplus C_3.
 \tag{E86.6}
\]
Here and below $C_m=\mathbb Z/m$. These are actual Artin
coordinates in the computed genus fields. In the coordinates
of E.76, the restriction and kernel maps for $B_2$ are
\[
 \begin{gathered}
 \Gamma_B=C_9\oplus C_3^3,\qquad
 \chi_B(x,y,z,t)=x+2y+z+t\pmod3,\\
 \operatorname{Cl}(B_2)\{3\}\xrightarrow{\sim}\ker\chi_B,
 \quad(s,u,v)\longmapsto(-2s-3u,u,v,-s-2u-v),\\
 (x,y,z,t)\longmapsto(x-3y-3z-3t,y,z).
 \end{gathered}
 \tag{E86.7}
\]
The two displayed arrows are inverse modulo the indicated
cyclic orders. For the lower genus field computed in the
certificate, the corresponding maps are
\[
 \begin{gathered}
 \Gamma_L=C_9\oplus C_3^2,\qquad
 \chi_L(x,y,z)=2(x+y+z)\pmod3,\\
 \operatorname{Cl}(L_2)\{3\}\xrightarrow{\sim}\ker\chi_L,
 \quad(a,b)\longmapsto(-2a,b,-a-b),\\
 (x,y,z)\longmapsto(x-3y-3z,y).
 \end{gathered}
 \tag{E86.8}
\]
In (E86.8), the ordinary classes of the ideals in (E86.4)
are exactly $a_2=(1,0)$ and $b_2=(0,1)$, and
$[\mathfrak q_{3,2}]=(3,0)=3a_2$.
Both computed genus fields are the full Hilbert three-class
fields of their respective upper fields.

For clarity, the ordinary norm and extension matrices are
also explicit. Columns are source coordinates, the upper
coordinates are $(s,u,v)$ from (E86.7), the $L_2$
coordinates are $(a_2,b_2)$, and the $B_1$ coordinates
are the certified ones of E.76:
\[
 \begin{array}{c|c|c}
 \text{extension}&N&j\\ \hline
 B_2/L_2&
 \begin{pmatrix}3&3&6\\2&0&0\end{pmatrix}&
 \begin{pmatrix}0&0\\1&0\\0&0\end{pmatrix}\\[3pt]
 B_2/B_1&
 \begin{pmatrix}1&1&0\\2&0&0\end{pmatrix}&
 \begin{pmatrix}0&6\\0&0\\0&0\end{pmatrix}.
 \end{array}
 \tag{E86.9}
\]
The subgroup generated by the six S-primes in $B_2$ is
exactly $(\operatorname{Cl}(B_2)\{3\})[3]$, of order $27$.
In the ordering of E.81 their classes are
\[
 [Q_1]=[Q_2]=[Q_3]=(6,1,0),\qquad
 [T_1]=(6,2,2),\quad[T_2]=(0,2,2),\quad[T_3]=(3,2,2).
 \tag{E86.10}
\]
Consequently the full principal S-divisor lattice is
\[
 \begin{split}
 L_S=\{w\in\mathbb Z^6:;&w_1+w_2+w_3\equiv0\pmod3,\quad
                w_4\equiv w_5\equiv w_6\pmod3\},\\
 [\mathbb Z^6:L_S]&=27,\qquad
 H_S=\begin{pmatrix}
 3&2&2&0&0&0\\0&1&0&0&0&0\\0&0&1&0&0&0\\
 0&0&0&3&0&1\\0&0&0&0&3&1\\0&0&0&0&0&1
 \end{pmatrix}.
 \end{split}
 \tag{E86.11}
\]
This evaluates the exact residual classes of E.84:
$\delta=w=0$, its group $W=0$, and all twenty-seven
literal ideals $J_{a,b,c}$ of E.81 are principal.
The complete Smith coordinate maps for (E86.7)--(E86.11)
and actual prime-ideal generators are retained.

There are exact evaluated unit-norm defects at every level.
Write $\mathscr U_d=\mathcal O_{B_d}[1/3,1/17]^\times$
and $\mathscr V_d=\mathcal O_{L_d}[1/3,1/17]^\times$.
For $r=1,2$, $n=3^r$, and
$K_{d,r}=\mathscr U_d\cap N_{B_{d+r}/B_d}B_{d+r}^\times$,
\[
 \begin{gathered}
 K_{d,r}/N\mathscr U_{d+r}=C_3,\qquad
 \mathscr U_d/N\mathscr U_{d+r}\simeq C_n^5\oplus C_3,
 \qquad [\mathscr U_d:N\mathscr U_{d+r}]=3^{5r+1},\\
 \operatorname{Ann}_{\mathbb Z_3[\operatorname{Gal}(B_{d+r}/B_d)]}
                  [\mathfrak p_{d+r}]_3=(3,\sigma-1).
 \end{gathered}
 \tag{E86.12}
\]
For the actual adjacent cyclic extension $L_{d+1}/L_d$, put
$K'_d=\mathscr V_d\cap N L_{d+1}^\times$. Then
\[
 [\mathscr V_d:K'_d]=3,\qquad
 K'_d/N\mathscr V_{d+1}=C_3^2,\qquad
 \mathscr V_d/N\mathscr V_{d+1}=C_3^3.
 \tag{E86.13}
\]
In (E86.12) the first quotient is identified with the
prime-$73$ class, and in (E86.13) with the actual pair
$(a_{d+1},b_{d+1})$. The identifying map is explicit:
given a field norm $u=Nx$, solve
\[
 \operatorname{div}_{\notin S}(x)=(\sigma-1)D;
 \qquad \omega(u)=[D]_3
       \quad\text{in the indicated S-class coordinates}.
 \tag{E86.14}
\]
Its kernel is exactly the S-unit norms. It is independent
of both choices, and is onto. Thus the groups in
(E86.12)--(E86.13) come with a membership map, not just
their orders. The elementary-divisor identifications of
the larger quotients allow choices of character lifts.

For the horizontal unramified extension $B_d/L_d$, the
complete S-unit norm condition is
\[
 N_{B_d/L_d}\mathscr U_d
 =\begin{cases}
 \{u\in\mathscr V_1:v_3(u)+2v_{17}(u)\equiv0\pmod3\},&d=1,\\
 \{u\in\mathscr V_d:v_{17}(u)\equiv0\pmod3\},&d\ge2.
 \end{cases}
 \tag{E86.15}
\]
The valuations are normalized at the unique primes of
$L_d$ above the indicated rational primes.

All these finite quotients, maps and zero criteria remain
exact after completion for multiplication by $3^m$.
Every nonempty completed unit-norm fiber is a translate
of the actual norm kernel. Its abstract compact group
and rank are
\[
 \begin{array}{c|c}
 \text{norm on S-units}&\text{kernel}\\ \hline
 B_{d+r}/B_d, r=1,2&
 C_{3^{\min(r,2)}}\oplus
       \mathbb Z_3^{\,3^{d+1}(3^r-1)}\\
 L_{d+1}/L_d&C_3\oplus\mathbb Z_3^{\,2\cdot3^d}\\
 B_d/L_d&C_3\oplus\mathbb Z_3^{\,2\cdot3^d+4}.
 \end{array}
 \tag{E86.16}
\]
The finite torsion is the literal kernel on roots of unity;
the free coordinates are choices of an integral basis of
the displayed norm kernel. Norms and extensions on the
S-class groups at arbitrary finite distances are the
compositions of (E86.3), with these same specified prime
bases. The assertions about cyclic unit defects in
(E86.12) use exactly $r=1,2$.

This evaluates all previously retained S-class
annihilators and capitulation maps in E.79. The actual
values of the remaining ordinary prime classes
$[\mathfrak q_{3,d}]$ for $d\ge3$, and their eventual
exponent $h$ in E.69, are not computed by (E86.3).
Their comparison is the actual norm/extension square
E76.11 and the exact S-class quotient
$\operatorname{Cl}(L_d)\{3\}\twoheadrightarrow C_d^S$,
whose kernel is generated by $[\mathfrak q_{3,d}]$.
The manuscript anchor remains LM~\S51.16,
fols.~733v--736v, scans~172--178, through the actual
positional norm maps of E.46, E.69 and E.72.
\end{statement}

\begin{proof}
\emph{The all-level automorphisms and their class maps.}
We first recall the exact field comparison behind the
horizontal arrows. The element $17$ has a cube root
in $\mathbb Q_3$, since $-17\in1+9\mathbb Z_3$ and
the cube map $1+3\mathbb Z_3\to1+9\mathbb Z_3$ is
bijective. Hence $\xi=\beta_1/a$ already belongs to
the completion of $L_1$ at three, and therefore to
every later such completion. At seventeen the
completion of $K$ has residue field $\mathbb F_{289}$,
which contains $\mu_9$; its Teichmuller lifts give
the corresponding local inclusion of $\xi$. Thus
the base change of $K(\xi)/K$ to $L_d$ is split
at these two primes. It is unramified elsewhere
because $K(\xi)/K$ is unramified away from three.
It has degree three: the unique prime over seventeen
is totally ramified in $L_d/K$, while it splits
into three primes in $K(\xi)/K$, excluding a global
inclusion. This proves the actual unramified,
S-split cubic comparison at every level, with
the embeddings fixed by (E86.1).

The adjacent $B$ extensions are cyclic cubic,
and the two-step extensions cyclic ninth-degree,
because $\mu_9\subset B_d$. They are totally
ramified at their six S-primes. At seventeen use
$Y^{3^r}-\beta_d$; at three use
$(Y-1)^{3^r}-\beta_d$ and the uniformizer
$\beta_d+1$ established in E.76. Their discriminants
have no other prime support. The adjacent $L$
extension is cyclic cubic, totally ramified at
its two S-primes. These observations justify every
cyclic and ramification hypothesis used below.

On $B_{d+1}$ define the literal substitutions
\[
 \begin{array}{c|cc}
      &\xi&\beta_{d+1}\\ \hline
 g_d&\xi&\xi\beta_{d+1}\\
 h_d&\zeta\xi&\beta_{d+1}
 \end{array}
 \quad
 g_d^9=h_d^3=1,\quad g_d^3=\sigma_d,\quad
 h_dg_dh_d^{-1}=g_d^4.
 \tag{E86.17}
\]
Here $\sigma_d$ is the specified generator of
$\operatorname{Gal}(B_{d+1}/B_d)$.
Both substitutions preserve $\xi^6+\xi^3+1=0$
and $\beta_{d+1}^{3^{d+1}}=17\xi^3$;
their inverse substitutions are powers of
themselves. Thus they are actual automorphisms.
They preserve $B_d$: $g_d(\beta_d)=\zeta\beta_d$
and $h_d(\beta_d)=\beta_d$. Both commute with
$\sigma_d$. Their restrictions satisfy
$h_d|_{B_d}=h_{d-1}$ and, for $d\ge2$,
$g_d|_{B_d}=g_{d-1}^3=\sigma_{d-1}$.
At $d=1$ the restriction is exactly the base
automorphism $g_F$ in E84.13--15. These are the
restriction maps, without an assertion that the
sequence $g_d$ itself is compatible.

E.79 proves that actual ideal norm induces
$(A_{d+1})_{\sigma_d}\simeq A_d$, and proves
$A_1=C_3$. Suppose the latter group at level $d$
has the specified generator. The coherent prime
$\mathfrak p_{d+1}$ maps to it, so its class
generates $A_{d+1}$ as a module over
$R=\mathbb Z_3[C_3]$. Here is the finite
Nakayama argument: quotient by the generated
module; the quotient $Q$ satisfies
$Q=(3,\sigma_d-1)Q$. Modulo three the augmentation
ideal is nilpotent, so iteration gives $Q=3Q$,
and finiteness gives $Q=0$.

Write this cyclic module as $R/I$. Since $R$ is
commutative, every endomorphism of $R/I$ is
multiplication by the image of one. Its endomorphism
ring is therefore commutative. The automorphisms
$g_d,h_d$ are $R$-linear, since they commute with
$\sigma_d$ and preserve S. Their induced
endomorphisms commute. Their actual commutator
in (E86.17) is $\sigma_d$, which consequently
acts identically. Thus $A_{d+1}=(A_{d+1})_{\sigma_d}
=C_3$, and the actual norm is an isomorphism.
Extension and norm satisfy
$jN=1+\sigma_d+\sigma_d^2=3$ on this group;
norm surjectivity therefore gives $j=0$.
This proves the $B$ part of (E86.3) by induction.
The same group of order three has trivial action
under every three-power order automorphism.
For the actual ninth-degree cyclic steps this
gives the stated annihilator $(3,\sigma-1)$;
norm and extension at longer distances are their
literal compositions.

The residue equalities $22^3=63$ and $63^3=22$
in $\mathbb F_{73}$ prove the compatibility of
the prime-$73$ chain at every level. These roots
have order eight, so the indicated residue maps
are defined at all levels. The defining polynomials
are etale at $73$, and their degree-one residue
maps extend to the maximal orders. Norm of each
prime is its predecessor. This also proves that
the induction uses the specified prime class,
not an unspecified cyclic generator.

\emph{The entire ordinary class group of $B_2$.}
Put $C=\operatorname{Cl}(B_2)\{3\}$ and let $B$
be its subgroup generated by the six S-primes.
Then $C/B=A_2=C_3$. The action of $g=g_1$ on
this quotient is trivial, because its order
divides nine. Hence $(g-1)C\subset B$.
E.84 proves the literal permutations
$g:Q_1\mapsto Q_3\mapsto Q_2\mapsto Q_1$,
$g(T_i)=T_i$, each S-class has order three,
and $[Q_1Q_2Q_3]=0$. Therefore
$1+g+g^2$ is zero on $B$. It follows that
\[
 (\sigma-1)C=(g^3-1)C
             =(1+g+g^2)(g-1)C=0.
 \tag{E86.18}
\]
The computed ambiguous order in E.76 is $81$,
so $|C|=81$ and $|B|=27$.

The genus map of E.81 has image of order $27$
on $B$, hence is injective there. Its zero
column difference gives $\delta=[Q_2]-[Q_1]=0$.
The already proved S-capitulation means
$[I]\in B$, for $I=P_{73}\mathcal O_{B_2}$.
Thus $w=[I]+[T_2]-[T_1]\in B$, and its zero
genus value gives $w=0$. E84.8 now proves
principality of all twenty-seven literal ideals.

To identify the full group and its actual map,
let $\mathcal H$ be the Hilbert three-class field
of $B_2$. The extension $\mathcal H/B_1$ is
Galois; its commutator subgroup is
$(\sigma-1)C$, because its quotient by $C$ is
cyclic. By (E86.18) it is abelian. Hence
$\mathcal H$ equals the maximal abelian-over-$B_1$
unramified-over-$B_2$ field computed in E.76.
Its group is $\Gamma_B=C_9\oplus C_3^3$.
The actual restriction character is computed
on prime Frobenii by
\[
 \overline{\beta_1}^{\,(N\mathfrak l-1)/3}
          =\overline\zeta^{\,\chi_B(\mathfrak l)}.
 \tag{E86.19}
\]
The four basis primes in the certificate have
norms $2^6,7^3,19^3,73$, respectively. Their
genus-coordinate matrix and character values are
\[
 V_B=\begin{pmatrix}3&3&6&5\\0&2&0&0\\
                    0&0&2&2\\1&0&2&2\end{pmatrix},
 \qquad (1,1,1,0).
 \tag{E86.20}
\]
The matrix is invertible modulo three, so these
primes generate $\Gamma_B$. Formula (E86.19)
therefore determines the character uniquely as
$(1,2,1,1)$; the program checks another $188$
actual prime Frobenii. Its full lattice kernel,
with integer Smith transformations and inverses,
is precisely (E86.7). Substituting the six
Artin columns of E.81 gives (E86.10), and gives
$[I]=(6,0,0)$. These classes generate the full
three-torsion subgroup of $C$, proving (E86.11)
in both directions. The full genus map is now
injective on $C$, so the earlier group $W$ is zero.

\emph{The horizontal S-class groups and all lower bases.}
For any $d$, let $\mathcal H_B^S$ be the S-split
Hilbert three-class field of $B_d$. It is Galois
over $L_d$. Its kernel over $B_d$ is $A_d=C_3$,
and conjugation has order dividing three, so
acts trivially on this kernel. An extension
of a cyclic group by a central cyclic group
is abelian. Thus $\mathcal H_B^S/L_d$ is
abelian of degree nine. It is everywhere
unramified and S-split.

The two field inclusions are useful here.
The extension $B_d/L_d$ is itself unramified,
abelian and S-split, so $B_d\subset\mathcal H_L^S$.
The extension $\mathcal H_L^S/B_d$ then has
the same properties, giving
$\mathcal H_L^S\subset\mathcal H_B^S$.
The abelian property just proved gives the
reverse inclusion. Hence
\[
 \mathcal H_L^S=\mathcal H_B^S,\qquad
                   |C_d^S|=9.
 \tag{E86.21}
\]
Restriction identifies actual norm $A_d\to C_d^S$
with the inclusion of the subgroup over $B_d$:
for a prime of residue degree $f$ it sends its
Frobenius to the $f$th power of the lower
Frobenius, exactly its ideal norm. Extension
in this common abelian class field is cubing.
Indeed a split lower prime extends to three
primes with the same Frobenius; an inert prime
has a single upper Frobenius equal to the cube.
Multiplication proves the assertion on all
ideal classes.

We determine its exponent without assuming an
upper unit basis. For an unramified S-split
cyclic cubic $E/F$, the invariant divisors
outside S are exactly extended divisors.
Hilbert~90 therefore gives the actual
isomorphism
\[
 \begin{gathered}
 \ker(j:\operatorname{Cl}_S(F)\to\operatorname{Cl}_S(E))
       \simeq H^1(\operatorname{Gal}(E/F),\mathcal O_{E,S}^\times),\\
 [D]\longmapsto\bigl(g\mapsto g(x)/x\bigr),
 \quad \operatorname{div}_{\notin S}(x)=jD.
 \end{gathered}
 \tag{E86.22}
\]
Changing $x$ by an S-unit changes a coboundary;
Hilbert~90 gives the inverse, and invariant
divisors give precisely the displayed kernel.
The S-unit Herbrand quotient in this extension
is $1/3$. All archimedean and S-place modules
are regular permutation modules, and the
product formula removes one trivial module.
The lattice proof of Milne VII Proposition~3.1
thus gives
\[
 |\ker j|=3[\mathcal O_{F,S}^\times:
                         N\mathcal O_{E,S}^\times].
 \tag{E86.23}
\]
Finite-index changes of logarithmic lattices
and finite torsion have Herbrand quotient one
\cite{E_MilneCFT}.

For $d\ge2$, norm from $B_d$ to $B_2$ preserves
each of its six S-valuation coefficients:
there is one prime above each and residue
degree one. Consequently every principal
S-divisor in $B_d$ satisfies the congruences
(E86.11), after applying norm. In particular
the total exponent of its three primes above
seventeen is divisible by three. In $B_d/L_d$
this total is the normalized valuation of the
norm at $\mathfrak q_{17,d}$. But $\beta_d$
is an S-unit of $L_d$ with that valuation one.
It is therefore excluded from the S-unit
norm image. Thus its index is at least three.
By (E86.21)--(E86.23) the index is at most
three. It is exactly three, the whole group
$C_d^S$ capitulates, and the necessary
valuation congruence is sufficient.
Since extension was cubing in (E86.21),
$C_d^S$ has exponent three. This proves
$C_d^S=C_3^2$ and (E86.15) for $d\ge2$.

At $d=1$, the certified class group of $L_1$
is $C_3^2$ and its primes at three and seventeen
are principal (E.57). Thus $C_1^S=C_3^2$;
cubing again gives zero extension, and
(E86.23) gives unit-norm index three.
The certified principal S-divisor lattice in
$B_1$ is E79.17, whose necessary norm congruence
is exactly $v_3+2v_{17}=0$ modulo three.
It has index three, proving (E86.15) at this
remaining level as well.

Norm on S-classes in the adjacent $L$ extension
is onto. One direct class-field proof is to
base change its lower S-split Hilbert field;
its intersection with the cyclic ramified
extension is the lower field. The resulting
unramified S-split extension upstairs has
restriction map on prime Frobenii equal to
ideal norm. Both S-class groups have order
nine, so norm is an isomorphism. It follows
that the cyclic action on the upper group is
trivial, and $jN=3$ gives zero extension.

At seven the residue map $\beta_d=-1$ is
defined at every level, with $\zeta=2$.
It is etale and degree one, and cubing sends
it to its predecessor. At seventy-three the
same assertion was already proved. The
lower calculation below gives the indicated
basis at $d=2$; norm isomorphisms propagate
it upward and downward. At $L_1$, the
certificate independently gives the columns
$(2,0)$ and $(2,1)$ in its certified class
basis, an invertible matrix modulo three.
For the extension adjoining $\xi$, the
Frobenius at $P_d$ sends $\xi$ to
$\xi^7=\zeta^2\xi$, while at $J_d$ it
sends $\xi$ to $\xi^{73}=\xi$.
Thus its Artin coordinates on $a_d,b_d$
are exactly $1,0$. The prime $\mathfrak p_d$
has norm $J_d$, proving every map in
(E86.3)--(E86.5).

\emph{The ordinary lower group and the actual matrix comparison.}
The kernel of the ordinary-to-S class quotient
for $L_2$ is generated by $\mathfrak q_{3,2}$,
since $\mathfrak q_{17,2}=(\beta_2)$.
E.72 proves that the former has exact order
three. Therefore $|\operatorname{Cl}(L_2)\{3\}|=27$.
For its lower genus field the retained program
uses only the certified degree-six base
$L_1$, with polynomial $X^6+17X^3+289$.
It computes the local norm units for the
actual cubic $Y^3-\beta_1$, using the forty
literal norms $N(1+(Y+1)^k)$.
The seven-coordinate local unit quotient
modulo cubes has a norm subspace of rank
six; local reciprocity proves that this is
its full norm kernel. At seventeen the
kernel is the residue cubes. Quotienting
the ray class group by these local norm
unit groups gives exactly the maximal
abelian-over-$L_1$ field unramified over
$L_2$, by the idelic map and converse
proved in E76.17. Its order over $L_1$
is $81$ and its Smith group is
$C_9\oplus C_3^2$. Its degree $27$ over
$L_2$ now equals the three-primary class number,
so it is the Hilbert three-class field.
This also proves that the cyclic action
of $\operatorname{Gal}(L_2/L_1)$ on the
ordinary class group is trivial.

The three prime Frobenii with genus-coordinate
columns
\[
 V_L=\begin{pmatrix}0&8&2\\0&0&1\\2&0&0\end{pmatrix},
 \qquad (1,1,0)
 \tag{E86.24}
\]
have norms $2^6,7,7$. Their matrix is
invertible modulo three. The literal
cubic-residue formula (E86.19), now in
$L_1$, gives the character $(2,2,2)$.
The program checks all $82$ sampled prime
Frobenii and retains the full integer
Smith maps, yielding (E86.8).

The degree-one primes $(7,\beta_2-3)$ and
$(73,\beta_2-63)$ have ordinary coordinates
$(8,1)$ and $(0,1)$, respectively.
The first differs in its residue label
from the coherent $P_2=(7,\beta_2+1)$.
There is an explicit comparison: with
$\sigma(\beta_2)=\zeta\beta_2$ and
$\zeta=2$ modulo this prime, application
of $\sigma^2$ changes the root $3$ to
$3\zeta^{-2}=6=-1$. Thus
\[
 \sigma^2(7,\beta_2-3)=(7,\beta_2+1).
 \tag{E86.25}
\]
Their classes agree by the just proved
trivial cyclic action, equivalently by
their identical lower genus Artin value
$(2,1,0)$. Taking their actual ideal ratio
with $J_2$ gives $(1,0)$, as asserted.
The wild prime has lower genus value
$(3,0,0)$, giving $[\mathfrak q_{3,2}]=3a_2$.

We give enough arithmetic data to reproduce
all four ordinary maps in (E86.9).
Three degree-one primes of $B_2$ above $73$
are specified by $(t,\beta_2)$ equal to
$(33,17),(31,17),(29,33)$, with
$t=a-\xi$ in $B_1$. Their upper class
coordinates, and their norms to $L_2$
and $B_1$, are the respective columns
\[
 M=\begin{pmatrix}2&5&8\\0&2&1\\2&1&2\end{pmatrix},
 \quad N_LM=\begin{pmatrix}0&0&3\\1&1&1\end{pmatrix},
 \quad N_BM=\begin{pmatrix}2&1&0\\1&1&1\end{pmatrix}.
 \tag{E86.26}
\]
Here $M$ is invertible modulo nine, with
inverse
$\left(\begin{smallmatrix}6&8&8\\1&3&8\\7&4&2\end{smallmatrix}\right)$.
Hence these are actual ideal generators
and determine both norm matrices. The
lower coordinates of their contractions
are obtained by the residue maps
$\beta_1=22,22,21$, respectively, followed
by (E86.8); all ideal HNFs are retained.

The actual $h$ and $g$ actions on the upper
class group, obtained by their action on
the retained ray ideals, are
\[
 h=\begin{pmatrix}1&0&6\\1&1&0\\1&0&1\end{pmatrix},
 \qquad g=\begin{pmatrix}4&0&0\\2&1&0\\0&0&1\end{pmatrix}
       \quad\text{on }C_9\oplus C_3^2.
 \tag{E86.27}
\]
They satisfy $1+h+h^2=0$. Thus extension
of $b_2=N[\mathfrak p_2]_3$ is zero.
The prime $P_2$ at seven is inert in
$B_2/L_2$, since $\mu_9$ is not in
$\mathbb F_7$. Its unique upper prime
contracts in $B_1$ to the degree-three
prime with $\beta_1=6$; its upper genus
value is $(3,2,0,2)$, hence its class
is $(0,2,0)$. Therefore extension of
$a_2=[J_2/P_2]$ is $(0,1,0)$.
This proves the first extension matrix.
The second is E84.2 transported through
(E86.7), using $[I]=(6,0,0)$.
The certificate verifies both identities
$Nj=3$ and $jN=\sum\sigma$ with their
actual source and target moduli.

\emph{The cyclic divisor sequences and completed fibers.}
For each cyclic step used in (E86.12) or
(E86.13), every ramified prime belongs
to S. On the divisor group outside S,
norm is residue degree times the sum
on each orbit. Consequently
\[
 \ker N=(\sigma-1)\operatorname{Div}_{\notin S}(E),
 \qquad
 \operatorname{Div}_{\notin S}(E)^{\sigma}
                 =j\operatorname{Div}_{\notin S}(F).
 \tag{E86.28}
\]
For $u=Nx$ an S-unit, these identities
produce the divisor $D$ in (E86.14).
Changing $D$ changes its class by an
extended base class, now zero by
(E86.3). Changing $x$ by a norm-one
element changes $D$ by a principal
divisor, by Hilbert~90. This proves
well-definedness. Every invariant
S-class has a representative $D$ for
which $(\sigma-1)D$ is principal
outside S; its generator has S-unit
norm. Thus the map is onto. The quotient of
invariant full S-classes by extended classes
is killed by the cyclic degree: the ideal
norm identity multiplies any invariant class
by that degree. Hence its prime-to-three part
is zero, justifying the three-primary projection.
Finally a zero class in this quotient means
$D=\operatorname{div}_{\notin S}(z)+jD_0$;
dividing $x$ by $\sigma(z)/z$ gives an
S-unit with the same norm. The reverse
inclusion uses $D=0$. This proves the
exact kernel and image of (E86.14).
It is the exact Hilbert~90 map E79.22,
with both class groups and the extension
map now evaluated.

For the $B$ steps E.79 proves the full
five-coordinate local Hilbert-symbol
image $C_n^5$ and the S-unit Herbrand
quotient $n^5$. The kernel of class
extension is now $C_3$, so (E86.22)
and that Herbrand quotient give
$[\mathscr U_d:N\mathscr U_{d+r}]=3n^5$.
Together with (E86.14) this proves
the exact sequence
$0\to C_3\to\mathscr U_d/N\mathscr U_{d+r}
\to C_n^5\to0$.
The middle group is killed by $n$,
since inclusion of a base S-unit has
norm its $n$th power. Choose preimages
of the five standard generators; their
orders divide $n$, so they split this
sequence as modules over $\mathbb Z/n$.
This proves (E86.12) with its stated
choice of elementary-divisor coordinates.

For the $L$ step the S-unit Herbrand
quotient is $3$: its two totally
ramified S-places contribute two
trivial modules, and the product
formula removes one. Class extension
has kernel $C_3^2$, so the unit-norm
index is $27$. Formula (E86.14) has
kernel quotient $C_3^2$, hence the
field-norm index is three. Equivalently,
the two actual local Hilbert symbols
have the full product-one image $C_3$.
The cyclic Hasse norm theorem identifies
their simultaneous kernel with $K'_d$
\cite[VIII, Theorem~3.1]{E_MilneCFT}.
The full S-unit quotient is killed by
three, so is $C_3^3$, proving (E86.13).

There are literal generator procedures
for the defect coordinates. For the
$B$ step first represent
$[\mathfrak p_{d+r}]_3$ by an actual ideal.
Write the exponent of the full
S-class group as $3^s m$ with $(3,m)=1$,
choose $e\equiv1\pmod{3^s}$ and
$e\equiv0\pmod m$, and take
$D=\mathfrak p_{d+r}^{\,e}$.
This ideal represents the desired primary
projection. Choose a generator outside S
for $(\sigma-1)D$;
its field norm represents $1$ under
$\omega$. For the $L$ step do the same
with the actual ideals $J_{d+1}/P_{d+1}$
and $J_{d+1}$, applying the same primary
projection. Different choices change
only the S-unit-norm coset, by the
proved kernel statement. These procedures
and (E86.14) specify all quotient maps
with no ambiguity from the prime-to-three
class components. Alternatively the ideal
representative can be chosen directly by
its specified three-primary S-class.

Every S-unit group here is finitely
generated. Norm image has finite index,
since it contains the appropriate powers
of the base group. Completion for
multiplication by $3^m$ therefore preserves
all these kernels, images and finite
quotients. The ranks of the completed
S-unit groups are
$3^{d+1}+5$ for $B_d$ and $3^d+1$
for $L_d$. Surjectivity after tensoring
with $\mathbb Q_3$ gives the free-rank
differences in (E86.16).

The three-power roots of unity of $B_d$
are exactly $\mu_9$: at seventeen any
root of unity of order prime to seventeen
injects into $\mathbb F_{289}^\times$,
whose three-part has order nine.
For $L_d$ they are exactly $\mu_3$,
since adjoining $\xi$ has degree three.
In the vertical $B$ steps the root-of-unity
norm is raising to $3^r$; in the vertical
$L$ step it is cubing. Horizontally
$N_{B_d/L_d}\xi=\xi^{1+4+7}=\zeta$.
These give exactly the torsion subgroups
shown in (E86.16). A finitely generated
module over $\mathbb Z_3$ is its finite
torsion plus a free module, so choosing
an integral basis of its torsion-free
quotient and lifts gives the asserted
compact coordinates for every actual
kernel. Translation by one norm preimage
gives, with its inverse, every nonempty
fiber.
\end{proof}
